% Options for packages loaded elsewhere
\PassOptionsToPackage{unicode}{hyperref}
\PassOptionsToPackage{hyphens}{url}
\documentclass[
]{article}
\usepackage{xcolor}
\usepackage[margin=1in]{geometry}
\usepackage{amsmath,amssymb}
\setcounter{secnumdepth}{-\maxdimen} % remove section numbering
\usepackage{iftex}
\ifPDFTeX
  \usepackage[T1]{fontenc}
  \usepackage[utf8]{inputenc}
  \usepackage{textcomp} % provide euro and other symbols
\else % if luatex or xetex
  \usepackage{unicode-math} % this also loads fontspec
  \defaultfontfeatures{Scale=MatchLowercase}
  \defaultfontfeatures[\rmfamily]{Ligatures=TeX,Scale=1}
\fi
\usepackage{lmodern}
\ifPDFTeX\else
  % xetex/luatex font selection
\fi
% Use upquote if available, for straight quotes in verbatim environments
\IfFileExists{upquote.sty}{\usepackage{upquote}}{}
\IfFileExists{microtype.sty}{% use microtype if available
  \usepackage[]{microtype}
  \UseMicrotypeSet[protrusion]{basicmath} % disable protrusion for tt fonts
}{}
\makeatletter
\@ifundefined{KOMAClassName}{% if non-KOMA class
  \IfFileExists{parskip.sty}{%
    \usepackage{parskip}
  }{% else
    \setlength{\parindent}{0pt}
    \setlength{\parskip}{6pt plus 2pt minus 1pt}}
}{% if KOMA class
  \KOMAoptions{parskip=half}}
\makeatother
\usepackage{color}
\usepackage{fancyvrb}
\newcommand{\VerbBar}{|}
\newcommand{\VERB}{\Verb[commandchars=\\\{\}]}
\DefineVerbatimEnvironment{Highlighting}{Verbatim}{commandchars=\\\{\}}
% Add ',fontsize=\small' for more characters per line
\usepackage{framed}
\definecolor{shadecolor}{RGB}{248,248,248}
\newenvironment{Shaded}{\begin{snugshade}}{\end{snugshade}}
\newcommand{\AlertTok}[1]{\textcolor[rgb]{0.94,0.16,0.16}{#1}}
\newcommand{\AnnotationTok}[1]{\textcolor[rgb]{0.56,0.35,0.01}{\textbf{\textit{#1}}}}
\newcommand{\AttributeTok}[1]{\textcolor[rgb]{0.13,0.29,0.53}{#1}}
\newcommand{\BaseNTok}[1]{\textcolor[rgb]{0.00,0.00,0.81}{#1}}
\newcommand{\BuiltInTok}[1]{#1}
\newcommand{\CharTok}[1]{\textcolor[rgb]{0.31,0.60,0.02}{#1}}
\newcommand{\CommentTok}[1]{\textcolor[rgb]{0.56,0.35,0.01}{\textit{#1}}}
\newcommand{\CommentVarTok}[1]{\textcolor[rgb]{0.56,0.35,0.01}{\textbf{\textit{#1}}}}
\newcommand{\ConstantTok}[1]{\textcolor[rgb]{0.56,0.35,0.01}{#1}}
\newcommand{\ControlFlowTok}[1]{\textcolor[rgb]{0.13,0.29,0.53}{\textbf{#1}}}
\newcommand{\DataTypeTok}[1]{\textcolor[rgb]{0.13,0.29,0.53}{#1}}
\newcommand{\DecValTok}[1]{\textcolor[rgb]{0.00,0.00,0.81}{#1}}
\newcommand{\DocumentationTok}[1]{\textcolor[rgb]{0.56,0.35,0.01}{\textbf{\textit{#1}}}}
\newcommand{\ErrorTok}[1]{\textcolor[rgb]{0.64,0.00,0.00}{\textbf{#1}}}
\newcommand{\ExtensionTok}[1]{#1}
\newcommand{\FloatTok}[1]{\textcolor[rgb]{0.00,0.00,0.81}{#1}}
\newcommand{\FunctionTok}[1]{\textcolor[rgb]{0.13,0.29,0.53}{\textbf{#1}}}
\newcommand{\ImportTok}[1]{#1}
\newcommand{\InformationTok}[1]{\textcolor[rgb]{0.56,0.35,0.01}{\textbf{\textit{#1}}}}
\newcommand{\KeywordTok}[1]{\textcolor[rgb]{0.13,0.29,0.53}{\textbf{#1}}}
\newcommand{\NormalTok}[1]{#1}
\newcommand{\OperatorTok}[1]{\textcolor[rgb]{0.81,0.36,0.00}{\textbf{#1}}}
\newcommand{\OtherTok}[1]{\textcolor[rgb]{0.56,0.35,0.01}{#1}}
\newcommand{\PreprocessorTok}[1]{\textcolor[rgb]{0.56,0.35,0.01}{\textit{#1}}}
\newcommand{\RegionMarkerTok}[1]{#1}
\newcommand{\SpecialCharTok}[1]{\textcolor[rgb]{0.81,0.36,0.00}{\textbf{#1}}}
\newcommand{\SpecialStringTok}[1]{\textcolor[rgb]{0.31,0.60,0.02}{#1}}
\newcommand{\StringTok}[1]{\textcolor[rgb]{0.31,0.60,0.02}{#1}}
\newcommand{\VariableTok}[1]{\textcolor[rgb]{0.00,0.00,0.00}{#1}}
\newcommand{\VerbatimStringTok}[1]{\textcolor[rgb]{0.31,0.60,0.02}{#1}}
\newcommand{\WarningTok}[1]{\textcolor[rgb]{0.56,0.35,0.01}{\textbf{\textit{#1}}}}
\usepackage{graphicx}
\makeatletter
\newsavebox\pandoc@box
\newcommand*\pandocbounded[1]{% scales image to fit in text height/width
  \sbox\pandoc@box{#1}%
  \Gscale@div\@tempa{\textheight}{\dimexpr\ht\pandoc@box+\dp\pandoc@box\relax}%
  \Gscale@div\@tempb{\linewidth}{\wd\pandoc@box}%
  \ifdim\@tempb\p@<\@tempa\p@\let\@tempa\@tempb\fi% select the smaller of both
  \ifdim\@tempa\p@<\p@\scalebox{\@tempa}{\usebox\pandoc@box}%
  \else\usebox{\pandoc@box}%
  \fi%
}
% Set default figure placement to htbp
\def\fps@figure{htbp}
\makeatother
\setlength{\emergencystretch}{3em} % prevent overfull lines
\providecommand{\tightlist}{%
  \setlength{\itemsep}{0pt}\setlength{\parskip}{0pt}}
\usepackage{bookmark}
\IfFileExists{xurl.sty}{\usepackage{xurl}}{} % add URL line breaks if available
\urlstyle{same}
\hypersetup{
  pdftitle={Pre-registration and statistical analysis plans},
  hidelinks,
  pdfcreator={LaTeX via pandoc}}

\title{Pre-registration and statistical analysis plans}
\author{}
\date{\vspace{-2.5em}}

\begin{document}
\maketitle

\textbf{Workflow lab:} Goal → Approach → Execution → Check → Report.

\subsection{Learning objectives}\label{learning-objectives}

\begin{itemize}
\tightlist
\item
  Distinguish a pre-registration from a statistical analysis plan.
\item
  Draft the core sections of each for a realistic study.
\item
  Recognise the forms of post-hoc flexibility that pre-registration is
  designed to prevent.
\end{itemize}

\subsection{Prerequisites}\label{prerequisites}

Earlier design and causal-inference material.

\subsection{Background}\label{background}

A pre-registration is a timestamped public record of a study's
hypotheses, design, and primary analysis. A statistical analysis plan
(SAP) is a more detailed, often confidential companion document that
specifies exactly how the data will be handled --- variable definitions,
handling of missing data, subgroup analyses, sensitivity analyses, and
reporting conventions. Between them, these two documents make the
difference between ``we planned this'' and ``we planned this, we can
prove it, and here is the file to show when it was signed.''

Flexibility during analysis --- \emph{researcher degrees of freedom} ---
is not fraud. It is usually well-meaning curiosity. But aggregated
across a field it produces a literature of spurious findings, and for
any one study it produces an analysis that will not replicate. A
pre-registration does not forbid curiosity. It separates
\emph{confirmatory} analyses (specified in advance) from
\emph{exploratory} ones, and requires each to be labelled when reported.

\subsection{Setup}\label{setup}

\begin{Shaded}
\begin{Highlighting}[]
\FunctionTok{library}\NormalTok{(tidyverse)}
\FunctionTok{set.seed}\NormalTok{(}\DecValTok{42}\NormalTok{)}
\FunctionTok{theme\_set}\NormalTok{(}\FunctionTok{theme\_minimal}\NormalTok{(}\AttributeTok{base\_size =} \DecValTok{12}\NormalTok{))}
\end{Highlighting}
\end{Shaded}

\subsection{1. Goal}\label{goal}

Produce a minimal-but-complete pre-registration template for a fictional
two-arm randomised trial, plus a SAP skeleton.

\subsection{2. Approach}\label{approach}

A pre-registration has four mandatory elements:

\begin{enumerate}
\def\labelenumi{\arabic{enumi}.}
\tightlist
\item
  \textbf{Question and hypotheses} --- exact form of H₀ and H₁ for each
  primary outcome.
\item
  \textbf{Design} --- eligibility, allocation, blinding, intervention,
  follow-up.
\item
  \textbf{Primary analysis} --- the single specification whose
  \emph{p}-value will be quoted as the headline result.
\item
  \textbf{Sample size justification} --- the calculation behind the
  target N.
\end{enumerate}

A SAP adds:

\begin{enumerate}
\def\labelenumi{\arabic{enumi}.}
\setcounter{enumi}{4}
\tightlist
\item
  \textbf{Secondary analyses} --- ordered and pre-specified.
\item
  \textbf{Handling of missing data} --- mechanism assumed, method used.
\item
  \textbf{Subgroup analyses} --- pre-specified, limited in number.
\item
  \textbf{Sensitivity analyses} --- especially for primary outcomes.
\end{enumerate}

\subsection{3. Execution --- template}\label{execution-template}

\begin{Shaded}
\begin{Highlighting}[]
\CommentTok{\# Pre{-}registration — Trial X}

\NormalTok{Protocol version }\FloatTok{1.0} \SpecialCharTok{|}\NormalTok{ Date }\DecValTok{2026{-}04{-}18} \SpecialCharTok{|}\NormalTok{ PI [NAME]}

\DocumentationTok{\#\# 1. Research question}
\NormalTok{Primary}\SpecialCharTok{:}\NormalTok{ Does [intervention] reduce [outcome] relative to [control]}
\ControlFlowTok{in}\NormalTok{ [population] over [time frame]?}

\DocumentationTok{\#\# 2. Hypotheses}
\NormalTok{H0}\SpecialCharTok{:}\NormalTok{ difference }\ControlFlowTok{in}\NormalTok{ [outcome] }\OtherTok{=} \FloatTok{0.}
\NormalTok{H1}\SpecialCharTok{:}\NormalTok{ difference }\ControlFlowTok{in}\NormalTok{ [outcome] }\SpecialCharTok{!=} \FloatTok{0.}\NormalTok{ Two}\SpecialCharTok{{-}}\NormalTok{sided, alpha }\OtherTok{=} \DecValTok{0}\NormalTok{.}\FloatTok{05.}

\DocumentationTok{\#\# 3. Design}
\NormalTok{Two}\SpecialCharTok{{-}}\NormalTok{arm, parallel}\SpecialCharTok{{-}}\NormalTok{group, double}\SpecialCharTok{{-}}\NormalTok{blind, placebo}\SpecialCharTok{{-}}\NormalTok{controlled trial.}
\DecValTok{1}\SpecialCharTok{:}\DecValTok{1}\NormalTok{ allocation, block }\FunctionTok{randomisation}\NormalTok{ (block size }\DecValTok{4}\NormalTok{), stratified by site.}

\DocumentationTok{\#\# 4. Primary analysis}
\NormalTok{ITT. Linear regression of outcome at follow}\SpecialCharTok{{-}}\NormalTok{up on arm, adjusted }\ControlFlowTok{for}
\NormalTok{baseline value and stratification factors. Primary estimand}\SpecialCharTok{:}
\NormalTok{adjusted mean difference with }\DecValTok{95}\NormalTok{\% CI.}

\DocumentationTok{\#\# 5. Sample size}
\NormalTok{Target n }\OtherTok{=} \DecValTok{200}\NormalTok{ per arm; power }\OtherTok{=} \FloatTok{0.80} \ControlFlowTok{for}\NormalTok{ d }\OtherTok{=} \FloatTok{0.3}\NormalTok{, alpha }\OtherTok{=} \DecValTok{0}\NormalTok{.}\FloatTok{05.}

\DocumentationTok{\#\# 6. Data handling}
\SpecialCharTok{{-}}\NormalTok{ Missing data}\SpecialCharTok{:}\NormalTok{ multiple imputation under }\FunctionTok{MAR}\NormalTok{ (}\AttributeTok{m =} \DecValTok{20}\NormalTok{).}
\SpecialCharTok{{-}}\NormalTok{ Outliers}\SpecialCharTok{:}\NormalTok{ retained }\ControlFlowTok{in}\NormalTok{ primary analysis.}
\SpecialCharTok{{-}}\NormalTok{ Adherence}\SpecialCharTok{:}\NormalTok{ per}\SpecialCharTok{{-}}\NormalTok{protocol analysis as sensitivity.}
\end{Highlighting}
\end{Shaded}

\subsection{4. Check}\label{check}

Three questions to ask of any pre-registration before posting:

\begin{enumerate}
\def\labelenumi{\arabic{enumi}.}
\tightlist
\item
  Could a reader reconstruct the primary analysis from the text without
  contacting you?
\item
  Are exploratory analyses labelled as such?
\item
  Is the target N justified by a calculation a reader can redo?
\end{enumerate}

\subsection{5. Report}\label{report}

\begin{quote}
A pre-registration for Trial X was deposited on the OSF on {[}date{]}
(DOI {[}doi{]}). The primary analysis, statistical model, and
sample-size justification are specified therein. Any deviation from the
pre-registered plan is reported in the Deviations section of the
manuscript with its rationale.
\end{quote}

\subsection{Common pitfalls}\label{common-pitfalls}

\begin{itemize}
\tightlist
\item
  Treating pre-registration as a checkbox without a concrete analysis
  plan.
\item
  Pre-registering two primary outcomes and quoting the winner.
\item
  Failing to report deviations from plan.
\end{itemize}

\subsection{Further reading}\label{further-reading}

\begin{itemize}
\tightlist
\item
  Nosek BA et al.~(2018). \emph{The preregistration revolution.}
\item
  AMRC / NIHR Statistical Analysis Plan templates.
\end{itemize}

\subsection{Session info}\label{session-info}

\begin{Shaded}
\begin{Highlighting}[]

\end{Highlighting}
\end{Shaded}

\subsection{See also --- chapter index}\label{see-also-chapter-index}

\begin{itemize}
\tightlist
\item
  \href{../index.Rmd}{Methods in this chapter}
\item
  \href{../../../ch00-find-your-method.Rmd}{Find your method (Chapter
  0)}
\end{itemize}

\end{document}
